\begin{ex}%[2H5C2-6]%[ID]
	Trong KG $Oxyz$, cho điểm $A(0 ; 3 ;-2)$. Xét đường thẳng $d$ thay đổi song song với $Oz$ và cách $Oz$ một khoảng bằng $2$. Khi khoảng cách từ $A$ đến $d$ nhỏ nhất, biết $d$ đi qua điểm $M(a;b;c)$ sao cho $a$, $b$, $c$ lập thành một cấp số cộng. Tính $a^2+b^2-c^2$.
	\shortans{$-12$}
	\loigiai{
		Vì $d$ song song với $Oz$ và cách $Oz$ một khoảng bằng $2$ nên $d$ thuộc mặt trụ trục $Oz$ và bán kính bằng $2$.\\
		Có $H(0 ; 0 ;-2)$ là hình chiếu vuông góc của $A(0 ; 3 ;-2)$ trên $Oz$.\\
		Có $\overrightarrow{HA}=(0 ; 3 ; 0) \Rightarrow H A=3$ nên $A$ nằm ngoài mặt trụ.\\
		Gọi $(P)$ là mặt phẳng qua $A$ và vuông góc với $Oz$.\\
		$M$ là hình chiếu vuông góc của $A$ trên $d$.\\
		Gọi $K$ là giao điểm của $AH$ và mặt trụ ($K$ nằm giữa $A$ và $H$).\\
		Dễ thấy $D(A ; d)=A M \geq A K$; $A K=A H-D(A ; d)=1$.\\
		Dấu bằng xảy ra khi và chỉ khi $M \equiv K$.\\
		Khi đó ta có $\overrightarrow{H K}=\dfrac{2}{3} \overrightarrow{H A} \Rightarrow K(0 ; 2 ;-2) \Rightarrow d\colon \heva{&x=0 \\ &y=2 \\ &z=-2+t} (t \in \mathbb{R})$.\\
		Vậy điểm $M \in d$ sẽ là $M(0;2;4)$ suy ra $a^2+b^2-c^2=-12$.
	}
\end{ex}

\begin{ex}%[2H5C2-7]%[ID]
	Trong KG $Oxyz$, đường thẳng $\Delta$ đi qua điểm $M(3 ; 1 ; 1)$, nằm trong mặt phẳng $(\alpha)\colon x+y-z-3=0$ và tạo với đường thẳng $d\colon \heva{&x=1 \\ &y=4+3 t \\ &z=-3-2 t}$ một góc nhỏ nhất. Biết điểm $N(a;b;2) \in d$, tính $\dfrac{4b}{a}$.
	\shortans{$-1{,}5$}
	\loigiai{
		\begin{itemize}
			\item \textbf{Cách 1.}\\
			Gọi $d'$ là hình chiếu vuông góc của $d$ lên $(\alpha)$, khi đó góc $\left(d, d'\right)$ là góc nhỏ nhất trong các góc tạo bởi $d$ với đường thẳng bất kỳ trong $(\alpha)$.\\
			$d$ có véc-tơ chỉ phương $\overrightarrow{u}_d=(0 ; 3 ;-2)$, $(\alpha)$ có véc-tơ pháp tuyến $\overrightarrow{n}_\alpha=(1 ; 1 ;-1)$.\\
			Gọi $\overrightarrow{n}_\beta$ là véc-tơ pháp tuyến của mặt phẳng $(\beta)$ tạo bởi $d$ và $d'$ thì $\overrightarrow{n}_\beta \perp \overrightarrow{u}_d$, $\overrightarrow{n}_\beta \perp \overrightarrow{n}_\alpha$ ta chọn $\overrightarrow{n}_\beta=\left[\overrightarrow{n}_\alpha, \overrightarrow{u}_d\right]=(1 ; 2 ; 3)$.\\
			Gọi $\overrightarrow{u}_{d'}$ là véc-tơ chỉ phương của $d'$ thì $\overrightarrow{u}_{d'} \perp \overrightarrow{n}_\alpha, \overrightarrow{u}_{d'} \perp \overrightarrow{n}_\beta$, ta chọn $\overrightarrow{u}_{d'}=\left[\overrightarrow{n}_\alpha, \overrightarrow{n}_\beta\right]=(5 ;-4 ; 1)$.\\
			Đường thẳng $\Delta$ song song hoặc trùng với $d'$ nên có véc-tơ chỉ phương $\overrightarrow{u}_{\Delta}=\overrightarrow{u}_{d'}=(5 ;-4 ; 1)$.\\
			Vậy PTTS của đường thẳng $\Delta$ là $\heva{&x=3+5t\\&y=1-4t\\&z=1+t.} (t \in \mathbb{R})$\\
			Do $N(a;b;2) \in \Delta \Leftrightarrow \heva{&a=3+5t\\&b=1-4t\\&2=1+t} \Leftrightarrow\heva{&a=8\\&b=-3\\&t=1.} $ nên $\dfrac{4b}{a}=-1{,}5$.
			\item \textbf{Cách 2.}\\
			$d$ có véc-tơ chỉ phương $\overrightarrow{u}_d=(0 ; 3 ;-2)$, $(\alpha)$ có véc-tơ pháp tuyến $\overrightarrow{n}_\alpha=(1 ; 1 ;-1)$.\\
			Giả sử $\Delta$ có véc-tơ chỉ phương $\overrightarrow{u}_{\Delta}=(a ; b ; c),\left(a^2+b^2+c^2>0\right)$.\\
			Do $\Delta \subset(\alpha)$ ta có $\overrightarrow{n}_\alpha \cdot \overrightarrow{u}_{\Delta}=0 \Leftrightarrow a+b-c=0 \Leftrightarrow c=a+b$.
			Gọi $\varphi$ là góc giữa $\Delta$ và $d$, $\varphi \in\left[0 ; \dfrac{\pi}{2}\right]$, khi đó
			$$
			\cos \varphi=\left|\cos \left(\overrightarrow{u}_{\Delta}, \overrightarrow{u}_d\right)\right|=\dfrac{|3 b-2 c|}{\sqrt{13} \cdot \sqrt{a^2+b^2+c^2}}=\dfrac{|b-2 a|}{\sqrt{13} \cdot \sqrt{2 a^2+2 b^2+2 a b}}
			$$
			Góc $\varphi$ nhỏ nhất khi và chỉ khi $\cos \varphi$ lớn nhất, ta xét các trường hợp
			\begin{itemize}
				\item \textbf{Trường hợp 1.} Nếu $a=0$ ta được $\cos \varphi=\dfrac{1}{\sqrt{26}}$.
				\item \textbf{Trường hợp 2.} Nếu $a \neq 0$ ta được $\cos \varphi=\dfrac{|t-2|}{\sqrt{26\left(t^2+t+1\right)}}$, $\left(t=\dfrac{b}{a}\right)$.\\
				Ta có $26 \cos ^2 \varphi=\dfrac{t^2-4 t+4}{t^2+t+1}$, đặt $f(t)=\dfrac{t^2-4 t+4}{t^2+t+1}, t \in \mathbb{R}$, có
				$$	f'(t)=\dfrac{5 t^2-6 t-8}{\left(t^2+t+1\right)^2},\ f'(t)=0 \Leftrightarrow\hoac{&t=2 \\&	t=-\dfrac{4}{5}.}$$
				Bảng biến thiên của hàm số $f(t)$ như sau
				\begin{center}
				\begin{tikzpicture}[>=stealth]
					\tkzTabInit[nocadre=false,lgt=1.2,espcl=2,deltacl=0.5]{$t$/1.2 ,$f'(t)$/.7,$f(t)$/2}
					{$-\infty$ , $-\dfrac{4}{5}$ , $2$ , $+\infty$}
					\tkzTabLine{ , + , $0$ , - , $0$ , + , }
					\tkzTabVar{-/$1$ , +/$\dfrac{28}{3}$ , -/$\dfrac{2}{3}$ , +/$1$}
				\end{tikzpicture}
				\end{center}
				Do $\cos \varphi \geq 0$ nên $\cos \varphi$ lớn nhất khi $f(t)$ lớn nhất, từ bảng biến thiên ta được $\max\limits_{t \in \mathbb{R}} f(t)=f\left(-\dfrac{4}{5}\right)$.\\
				Khi đó $\dfrac{b}{a}=-\dfrac{4}{5} \Leftrightarrow 5 b=-4 a$, chọn $a=5 \Rightarrow b=-4$ và ta được $c=1$.\\
				Đường thẳng $\Delta$ có véc-tơ chỉ phương $\overrightarrow{u}_{\Delta}=(5 ;-4 ; 1)$.\\
				Vậy PTTS của đường thẳng $\Delta$ là $\heva{&x=3+5t\\&y=1-4t\\&z=1+t.} (t \in \mathbb{R})$\\
				Do $N(a;b;2) \in \Delta \Leftrightarrow \heva{&a=3+5t\\&b=1-4t\\&2=1+t} \Leftrightarrow\heva{&a=8\\&b=-3\\&t=1} $ nên $\dfrac{4b}{a}=-1{,}5$.
			\end{itemize}
		\end{itemize}
	}
\end{ex}

\begin{ex}%[2H5C2-3]
	Trong không gian với hệ tọa độ $O x y z$, cho đường thẳng $(d)\colon \heva{&x=-1+2 t \\ &y=1-t \\ &z=2 t}$ và hai điểm $A(1 ; 5 ; 0)$, $B(3 ; 3 ; 6)$. Gọi $M(a ; b ; c)$ là điểm trên $(d)$ sao cho chu vi tam giác $M A B$ đạt giá trị nhỏ nhất. Tính $P=a+b+c$.
	\shortans{$3$}
	\loigiai{
			Ta có $M \in d \Rightarrow M(-1+2 t ; 1-t ; 2 t)$.\\
			Chu vi tam giác $M A B$ là $A M+B M+A B$.\\
			Vì $A B$ là hằng số nên chu vi nhỏ nhất khi $A M+B M$ nhỏ nhất.\\
			Ta có $\overrightarrow{A M}(2 t-2 ;-t-4 ; 2 t)$, $\overrightarrow{B M}(2 t-4 ;-t-2 ; 2 t-6)$.
			$$
			A M+B M=\sqrt{9 t^2+20}+\sqrt{9 t^2-36 t+56}=\sqrt{(3 t)^2+(2 \sqrt{5})^2}+\sqrt{(6-3 t)^2+(2 \sqrt{5})^2}
			$$
			Đặt $\overrightarrow{u}=(3 t ; 2 \sqrt{5})$, $\overrightarrow{v}=(6-3 t ; 2 \sqrt{5}) \Rightarrow \overrightarrow{u}+\overrightarrow{v}=(6 ; 4 \sqrt{5})$.\\
			Áp dụng bất đẳng thức véc-tơ  $|\overrightarrow{u}|+|\overrightarrow{v}| \geq|\overrightarrow{u}+\overrightarrow{v}|$.\\
			Dấu bằng xảy ra khi $\overrightarrow{u}$, $\overrightarrow{v}$ cùng hướng.\\
			Ta có $A M+B M=|\overrightarrow{u}|+|\overrightarrow{v}| \geq|\overrightarrow{u}+\overrightarrow{v}|=\sqrt{6^2+(4 \sqrt{5})^2}=2 \sqrt{29}$.\\
			Do đó $A M+B M$ nhỏ nhất khi tồn tại số $k$ dương sao cho $$\overrightarrow{u}=k \overrightarrow{v} \Leftrightarrow\heva{&3 t=k(6-3 t) \\ &2 \sqrt{5}=2 \sqrt{5} k}\Leftrightarrow\heva{&t=1 \\ &k=1.}$$
			Khi đó $M(1 ; 0 ; 2)$ nên $P=a+b+c=1+0+2=3$.
	}
\end{ex}

\begin{ex}%[2H5C2-5]%[ID]
	Trong không gian với hệ tọa độ $O x y z$, cho tứ diện $A B C D$ có $A(-1 ; 1 ; 6)$, $B(-3 ;-2 ;-4)$, $C(1 ; 2 ;-1)$, $D(2 ;-2 ; 0)$. Điểm $M(a ; b ; c)$ thuộc đường thẳng $C D$ sao cho tam giác $A B M$ có chu vi nhỏ nhất. Tính $T=a+b+c$.
	\shortans{$1$}
	\loigiai{
			Gọi $C_{A B M}$ là chu vi của tam giác $A B M$.\\
			$ \overrightarrow{A B}=(-2 ;-3 ;-10) \Rightarrow A B=\sqrt{113} $.\\
			$ \overrightarrow{A B}=(-2 ;-3 ;-10)$, $\overrightarrow{C D}=(1 ;-4 ; 1) \Rightarrow \overrightarrow{A B} \cdot \overrightarrow{C D}=-2+12-10=0 \Rightarrow A B \perp C D$.
			Gọi $(P)$ là mặt phẳng chứa đường thẳng $A B$ và vuông góc với đường thẳng $C D$.\\ $H$ là giao điểm của $(P)$ và đường thẳng $C D$.\\
			Phương trình mặt phẳng $(P)$ qua $A(-1 ; 1 ; 6)$ có véc-tơ pháp tuyến $\overrightarrow{C D}=(1 ;-4 ; 1)$ là
			$$	x-4 y+z-1=0.$$
			PTĐT $C D\colon \heva{&x=1+t \\ &y=2-4 t \\ &z=-1+t.}$\\
			$ H \in C D \Leftrightarrow H(1+t ; 2-4 t ;-1+t)$.\\
			$H \in(P) \Leftrightarrow 1+t-4(2-4 t)-1+t-1=0 \Leftrightarrow t=\dfrac{1}{2} \Leftrightarrow H\left(\dfrac{3}{2} ; 0 ;-\dfrac{1}{2}\right)$.\\
			Với mọi $ M \in C D$, ta có $\heva{&A M \geq A H \\ &B M \geq B H} \Rightarrow A M+B M \geq A H+B H$.\\
			$$	\mathrm{C}_{A B M}=A B+A M+B M \geq \sqrt{113}+A H+B H,\ \forall M \in C D \text {. }
			$$
			Suy ra $\min \mathrm{C}_{A B M}=\sqrt{113}+A H+B H$, đạt được $M \equiv H \Leftrightarrow M\left(\dfrac{3}{2} ; 0 ;-\dfrac{1}{2}\right)$.\\
			Vậy $T=a+b+c=1$.
	}
\end{ex}
\Closesolutionfile{ans}
\indapan{6}{ans/ans-C5B3CD6-KQ}